\begin{frame}{Презентация к защите: \ProgDefenceSlides\ слайдов при лимите \ProgDefenceLimit}
\SlideTableTitle{Изменения по рекомендациям методических книг}
{\ProgDense
\begin{tabularx}{\linewidth}{|L{0.33\linewidth}|L{0.25\linewidth}|Y|}
\hline Рекомендация & Источник & Изменение \\\hline
В заголовке --- главная мысль слайда; меньше --- значит лучше & Муромцева, с. 24--25 & Заголовки-утверждения, не более трёх пунктов \\\hline
Цель и гипотезу приводить дословно & Новиков, «Подготовка к защите» & Формулировки взяты из введения \\\hline
Схема лучше таблицы, таблица лучше текста; ритм изложения & Муромцева, с. 24--29 & Задачи показаны схемой; схемы чередуются с текстом \\\hline
Тезис, аргументы, демонстрация & Непряхин, гл. 2, с. 21--31 & Контроли показаны как исключаемые объяснения \\\hline
Регламент 10--15 минут, ответы готовить заранее & Новиков, «Подготовка к защите» & Два запасных слайда для ответов на вопросы \\\hline
\end{tabularx}\par}
\SlideGap
Добавлены объект и предмет, предполагаемая новизна, положения, значимость и апробация. Шрифты и сетка сохранены по согласованному образцу.
\end{frame}
